\chapter{The Logical Topology and Compactness}
\label{sec:topology}
% ===================================================================

\section{The Logical Topology}

\begin{definition}[Logical Topology]
\label{def:log-top}
Let $\GSLT$ be a GSLT with context-decorated HML $\HML(\ctx)$.
For each formula $\phi \in \HML(\ctx)$, define the
\emph{extension} of $\phi$:
\[
  \llb\phi\rrb \;=\; \{ [P] \in \terms(\GSLT)/{\bisim} \mid P \sat \phi \}.
\]
The \emph{logical topology} $\mathcal{O}_{\HML}$ on
$\terms(\GSLT)/{\bisim}$ is the topology generated by the
extensions $\{ \llb\phi\rrb \mid \phi \in \HML(\ctx) \}$ as a
subbasis.
\end{definition}

\begin{remark}
This topology is the standard \emph{Scott topology} on the
bisimulation quotient, familiar from domain theory and Stone duality.
The resulting topological space is a \emph{spectral space} --- sober,
$T_0$, and quasi-compact --- when the GSLT is well-behaved.
The logical metric $d_{\HML}$ from Part~I is compatible
with this topology: the metric balls are opens in $\mathcal{O}_{\HML}$.
\end{remark}

\section{Population Subspaces}

\begin{definition}[Latent Subspace]
Given a population $\pop \subseteq \terms(\GSLT)/{\bisim}$, the
\emph{latent topological subspace} of $\pop$ is the subspace
\[
  X_\pop \;=\; \overline{\pop} \;\subseteq\; \terms(\GSLT)/{\bisim}
\]
with the subspace topology inherited from $\mathcal{O}_{\HML}$.
The closure is taken in the logical topology.
\end{definition}

\begin{remark}
The latent subspace is ``latent'' in the sense that it is not
constructed deliberately by the agents.
It is the topological shadow cast by the population's collective
behavior --- the smallest closed set in the logical topology that
contains all the agents.
An agent inside $\pop$ cannot directly observe $X_\pop$;
it can only probe it indirectly through experiments.
\end{remark}

\section{Compactness as Consensus}

\begin{definition}[Compact Population]
A population $\pop$ is \emph{compact} if its latent subspace
$X_\pop$ is compact in the logical topology $\mathcal{O}_{\HML}$:
every open cover of $X_\pop$ by extensions of HML formulae has a
finite subcover.
\end{definition}

\begin{theorem}[Compactness Implies Consensus]
\label{thm:compact-consensus}
Let $\pop$ be a compact population in an interactive GSLT $\GSLT$.
Then $\pop$ supports consensus: there is no infinite oscillation
on any Boolean question expressible in $\HML(\ctx)$.

More precisely: for any formula $\phi \in \HML(\ctx)$, if the
population is divided into two communities indefinitely alternating
between $\phi$ and $\neg\phi$, then the alternation must terminate
in finite time.
\end{theorem}

\begin{proof}[Proof sketch]
Suppose for contradiction that a sub-population oscillates
indefinitely between $\phi$ and $\neg\phi$.
Define formulae:
\[
  \phi_1 = \langle 0 \rangle\top, \quad
  \phi_2 = \langle 1 \rangle\langle 0 \rangle\top, \quad
  \phi_3 = \langle 0 \rangle\langle 1 \rangle\langle 0 \rangle\top,
  \quad \ldots
\]
where $0$ and $1$ label the two communities' states.
Each $\phi_n$ is satisfied by agents that have completed at least
$n$ oscillation steps.
The family $\{ \llb \phi_n \rrb \}_{n \in \NN}$ is an open cover
of $X_\pop$ with no finite subcover: for any finite $N$,
agents that have oscillated more than $N$ times are not covered
by $\bigcup_{n \leq N} \llb \phi_n \rrb$.
This contradicts compactness of $X_\pop$.
Therefore the oscillation must terminate.
\end{proof}

\begin{remark}[The Proof Delivers the Protocol]
Theorem~\ref{thm:compact-consensus} is an existence result:
it says consensus \emph{must} be reached, without specifying how.
However, the proof of compactness for a \emph{specific} population
subspace is constructive: it identifies the finite subcover, which
encodes the maximum depth of oscillation and therefore the
convergence time.
The proof of compactness is the consensus protocol.
This is the precise sense in which the topology does computational
work.
\end{remark}

\begin{remark}[Non-Compact Populations]
A non-compact population may contain irresolvable oscillations ---
communities that never agree.
Non-compactness is the topological signature of
\emph{fundamental disagreement}: not a disagreement that more
communication or better arguments could resolve, but one that is
structurally permanent.
In the context of the massive population, non-compactness at a
given scale means the cluster at that scale cannot function as a
coherent agent.
\end{remark}

% ===================================================================
